\section{Conclusion}\label{sec:conclusion}
We presented AutoML-IoT, a controlled evaluation pipeline that studies how reliably AutoML selects anomaly detection models for IoT systems. While practitioners tend to compare frameworks under a convenient random split, AutoML-IoT audits the data, constructs leakage-aware protocols, and compares the frameworks under identical conditions. We performed an audit of a vibration dataset and a traffic dataset to validate the existence of evaluation bias. In a controlled comparison, we evaluated Auto-sklearn and H2O AutoML against default baselines, and we measured detection quality, deployment cost, and feature reliance. We found that the evaluation protocol and the data influence the results more than the choice of framework. Label ambiguity bounds the attainable accuracy on the traffic data, and class balancing did not improve the macro-averaged F1 score. The frameworks differ mainly in the size and latency of the selected models. The insights gained from our analysis are useful for evaluating AutoML on further IoT datasets and for selecting models that fit the limits of IoT gateways.

\section*{Acknowledgments}
\todo{Funding and acknowledgments to be completed by the authors.}
